\section{Optimal Network Reconfiguration}
\label{sec:onr}

\subsection{Problem Form and Runtime Scope}

This chapter documents the current ONR mixed-integer model for radial feeder
reconfiguration.  The runtime problem is a LinDistFlow-style MILP with binary
branch-status variables, feeder-flow variables, squared-voltage variables, and
connectivity-flow variables, solved through the native branch-and-cut engine.
The formulation is written over solver-visible canonical network components,
matching the power-flow and public OPF convention in
Sections~\ref{sec:powerflow} and~\ref{sec:opf}.  The binary switching decision
is therefore a canonical line or edge status, not a direct physical command on
every circuit breaker or switch object.

\subsection{Code Map}

The current implementation is concentrated in:
\begin{itemize}
  \item \texttt{src/network\_reconfiguration/topology\_analysis.cpp} for ONR model build, solve, and post-processing,
  \item \texttt{src/network\_reconfiguration/topology\_reconfiguration.cpp} for fault-triggered reconfiguration,
  \item the forwarded MIPSolvers engine for native branch-and-cut, simplex, and
        external HiGHS/Gurobi adapters.
\end{itemize}

\subsection{Pseudocode Map}

The code-aligned workflow summary lives in two places:
\begin{itemize}
  \item Algorithm A22 in Section~\ref{sec:algo-onr} for the application-level ONR solve flow,
  \item Section~\ref{sec:engine-solver-system} for the solver-kernel mechanics used underneath the ONR MILP.
\end{itemize}

\subsection{Current Status}

The ONR route is a live MILP application using the native branch-and-cut stack.
Its main current boundary is modeling scope: projected canonical AC components
are supported, but higher-level hybrid overlays such as VPP, microgrid, and
mobile-storage behavior are not yet injected directly into the ONR MILP.

This chapter distinguishes two runtime paths.  The steady ONR path in
\texttt{topology\_analysis.cpp} solves radial feeder reconfiguration for loss
reduction.  The fault-triggered path in \texttt{topology\_reconfiguration.cpp}
builds a restoration MILP with failed-line availability, source limits, load
pickup, and optional PF-style constraints.  They share graph and MILP concepts
but have different data assembly and fallback behavior.

\subsection{Canonical Model Space}

The ONR solver exposes both \texttt{ACSystem} and \texttt{HybridPowerSystem}
entry points.  The \texttt{HybridPowerSystem} overload first projects the rich
model through \texttt{project\_to\_canonical\_models()}, so the steady ONR MILP
is assembled from \(\widehat{\mathcal{S}}=\Pi_{\qual{canon}}(\mathcal{S})\).
After this projection it calls the AC-only builder on
\texttt{projected.ac}.  The steady path can therefore see
\begin{itemize}
  \item direct AC buses and branches,
  \item direct AC loads and generators,
  \item direct static generators, renewable generators, PV systems, storage, and charging stations,
  \item projected transformers, switches, flexible loads, asymmetric loads, chargers, and optional three-phase reductions,
  \item still no automatic injection of \texttt{VPP}, \texttt{Microgrid}, or \texttt{MobileStorage} overlays into the ONR MILP.
\end{itemize}
It does not include projected DC branches, DC sources, VSC couplings, or DCDC
devices in the steady loss-minimizing ONR model.

The fault-triggered path, \texttt{run\_topology\_reconfiguration()}, also starts
from \(\widehat{\mathcal{S}}\), but it keeps a unified AC/DC/VSC edge list:
\begin{align}
\mathcal{E}_{\qual{hyb}} =
\mathcal{E}_{ac}^{\qual{canon}}
\cup \mathcal{E}_{dc}^{\qual{canon}}
\cup \mathcal{E}_{vsc}^{\qual{canon}} .
\end{align}
Its binary \(\beta_e\) variables are canonical edge-status decisions over this
hybrid list.  Structured \texttt{BranchRef\{category,index\}} results preserve
whether an edge came from an AC branch, DC branch, or VSC coupling; the legacy
integer vectors remain ambiguous in hybrid cases.

\subsection{Decision Variables}

For an AC network with \(m=|\mathcal{E}_{ac}|\) branches and
\(n=|\mathcal{N}_{ac}|\) buses, the MILP uses the variable blocks
\begin{align}
x = \begin{bmatrix}
\alpha & P & Q & v & f & t
\end{bmatrix}^{\top},
\end{align}
where
\begin{align}
\alpha_b &\in \{0,1\} && \text{branch status}, \\
P_b,Q_b &\in \mathbb{R} && \text{branch active/reactive flow}, \\
v_i &\in \mathbb{R} && \text{squared bus voltage magnitude}, \\
f_b &\in \mathbb{R} && \text{single-commodity connectivity flow}, \\
t_b &\ge 0 && \text{auxiliary variable for } |P_b|.
\end{align}

For the hybrid fault-triggered path, the analogous topology status is
\(\beta_e\in\{0,1\}\) on \(e\in\mathcal{E}_{\qual{hyb}}\).  The model still
optimizes line or edge status; physical device status is outside the MILP
variable vector.

\subsection{Switching Candidates and Device Realization}

The current ONR formulation treats switchability as a property of canonical
branches or edges:
\begin{align}
\alpha_b &\in \{0,1\}, && b\in\mathcal{K}_{ac},\\
\beta_e  &\in \{0,1\}, && e\in\mathcal{K}_{\qual{hyb}},
\end{align}
where \(\mathcal{K}\) is supplied by the caller or inferred from normally-open
ties.  In the steady path, \texttt{ONROptions::switchable\_branch\_ids} contains
stable \texttt{ACBranch::index} values; an empty list makes every canonical AC
branch a candidate.  Non-candidates are fixed to their current
\texttt{in\_service} state, and structural bridges are fixed closed.  In the
fault-triggered path, \texttt{TopoReconfOptions::switchable\_branches} should be
used for new hybrid code because it carries the edge category as well as the
component index.

This is not yet the same as a complete circuit-breaker and switch operation
model.  The implemented MILP decides whether a canonical branch or edge is
energized.  It does not introduce separate binary variables for each physical
\texttt{Switch}, \texttt{CircuitBreaker}, or \texttt{DCCircuitBreaker}, and it
does not yet derive \(\mathcal{K}\) automatically from the one-side or
double-side breaker layout of each protected line.

The projection layer provides partial provenance.  \texttt{BranchExpandMap}
records AC branches synthesized from transformers and AC switchgear, but both
\texttt{Switch} and \texttt{CircuitBreaker} entries currently use the same
\texttt{BranchOriginType::Switch} tag.  DC circuit breakers are represented as
\texttt{DC\_Switch} graph edges rather than entries in this AC branch expansion
map.  Closed ideal AC switches and breakers can also be contracted by
\texttt{merge\_zero\_impedance\_buses()}, which is correct for PF/OPF
canonicalization but can erase the direct one-to-one edge needed for a switching
operation report.

The intended ONR contract should therefore be:
\begin{align}
z_\ell &= \Phi_\ell(u_{\ell,1},u_{\ell,2},\ldots),
&& z_\ell\in\{0,1\},
\end{align}
where \(z_\ell\) is the canonical line-status decision optimized by ONR and
\(u_{\ell,k}\) are the physical commands for the switches or circuit breakers
that can realize that line state.  Examples include
\begin{align}
z_\ell &= u_{\ell,\qual{sw}}, && \text{direct switch or tie line},\\
z_\ell &= u_{\ell,\qual{cb}}, && \text{single-side breaker control},\\
z_\ell &= u_{\ell,\qual{from}}\land u_{\ell,\qual{to}},
&& \text{double-side breaker control}.
\end{align}
The current code implements \(z_\ell\) or \(\beta_e\), while the
line-to-device map \(\Phi_\ell\) is only implicit in component metadata such as
\texttt{element\_type}, \texttt{element\_id}, terminal buses, and projection
provenance.

\subsection{Net-Injection Model}

The ONR builder computes per-bus net injections directly from the AC component
tables as
\begin{align}
P_i^{net} &= \frac{P_i^{load} + P_i^{cs} - P_i^{gen} - P_i^{sg} - P_i^{ren} - P_i^{pv} - P_i^{stor}}{S_{base}}, \\
Q_i^{net} &= \frac{Q_i^{load} + Q_i^{cs} - Q_i^{gen} - Q_i^{sg} - Q_i^{ren} - Q_i^{pv} - Q_i^{stor}}{S_{base}}.
\end{align}
Storage discharge is therefore treated as generation, and charging stations are
additional fixed loads.

\subsection{LinDistFlow MILP}

\paragraph{Objective.}
The objective minimizes a linearized active-loss proxy:
\begin{align}
\min_x \sum_{b=1}^{m} r_b t_b.
\end{align}

\paragraph{Tree cardinality.}
\begin{align}
\sum_{b=1}^{m} \alpha_b = n-1.
\end{align}

\paragraph{Commodity-flow connectivity.}
The solver uses a single-commodity-flow spanning-tree formulation.  For the root
bus,
\begin{align}
\sum_{b\in\delta^{-}(root)} f_b - \sum_{b\in\delta^{+}(root)} f_b = -(n-1),
\end{align}
and for every non-root bus,
\begin{align}
\sum_{b\in\delta^{-}(i)} f_b - \sum_{b\in\delta^{+}(i)} f_b = 1.
\end{align}
Capacity of \(f_b\) is tied to \(\alpha_b\) using big-\(M\) logic.

\paragraph{Power-balance constraints.}
For each non-root bus,
\begin{align}
\sum_{b\in\delta^{-}(i)} P_b - \sum_{b\in\delta^{+}(i)} P_b &= P_i^{net}, \\
\sum_{b\in\delta^{-}(i)} Q_b - \sum_{b\in\delta^{+}(i)} Q_b &= Q_i^{net}.
\end{align}

\paragraph{Voltage-drop constraints.}
The current implementation uses the LinDistFlow relation relaxed by big-\(M\):
\begin{align}
v_j - v_i + 2r_b P_b + 2x_b Q_b &\le M(1-\alpha_b), \\
-(v_j - v_i + 2r_b P_b + 2x_b Q_b) &\le M(1-\alpha_b).
\end{align}

\paragraph{Thermal and absolute-value constraints.}
\begin{align}
|P_b| &\le P_b^{max}\alpha_b, \\
|Q_b| &\le Q_b^{max}\alpha_b, \\
P_b &\le t_b, \\
-P_b &\le t_b.
\end{align}

\subsection{Branch-and-Cut Configuration}

The ONR path solves the MILP through the native branch-and-cut engine with the
following current tuning:
\begin{itemize}
  \item cut family: MIR cuts,
  \item root cut rounds: 5,
  \item cuts per round: 20,
  \item branching: pseudocost,
  \item node selection: hybrid,
  \item LP node solver: simplex enabled,
  \item feasibility pump: enabled.
\end{itemize}

\subsection{Fault-Triggered Reconfiguration Path}

\texttt{run\_topology\_reconfiguration()} builds its source set on the unified
AC/DC bus index.  AC generators, static generators, renewables, PV systems,
storage units, \texttt{ExternalGrid} rows, and the final SLACK-bus fallback map
to AC source nodes.  DC static generators, DC-side static generators, DC
storage, DC PV arrays, and \texttt{DC\_V} buses map to DC source nodes.  A network
with no available AC or DC source returns infeasible before constructing root
source variables.

Load demand is additive.  Active/reactive demand from explicit
\texttt{ACSystem::loads} is accumulated first, and direct
\texttt{ACBus::pd\_mw}/\texttt{qd\_mvar} demand is then added at the same bus.
DC load rows are also mapped onto the hybrid restoration graph when the DC bus
is present.  The implementation does not treat bus demand and load records as
mutually exclusive encodings.

Switchable restoration edges are category-aware in hybrid systems.  New code
should populate \texttt{TopoReconfOptions::switchable\_branches} with
\texttt{BranchRef\{category,index\}} values so AC branches, DC branches, and
VSC couplings that share the same numeric \texttt{index} remain distinct.  The
legacy \texttt{switchable\_branch\_ids} vector is retained for compatibility but
is interpreted only as AC \texttt{ACBranch::index} values.  With no explicit
switchable list, the implementation keeps the historical automatic detection of
normally-open AC/DC tie branches.

If the heuristic warm start fails to produce a feasible restoration plan, the
code builds the full MILP fallback.  It tries HiGHS when the adapter is
available and retries native B\&C if HiGHS is unavailable, returns failure, or
returns a malformed/wrong-sized solution vector.  The stale behavior where a
missing heuristic plan directly implied infeasibility is no longer the current
runtime path.  Solver certificates use exact normalized optimality statuses plus
the MIP gap, avoiding substring matches such as ``not optimal''.  Runtime trace
output in this path and the native branch-and-cut fallback honors
\texttt{verbose=false}, keeping batch and CI runs quiet unless verbose logging is
requested.

\subsection{Postprocessing and Verification}

The steady ONR path in \texttt{solve\_optimal\_reconfiguration()} still applies
the returned AC topology to an \texttt{ACSystem} and runs Newton PF as a
verification step.  The fault-triggered hybrid path
\texttt{run\_topology\_reconfiguration()} instead performs algebraic post-solve
checks on equality rows, inequality rows, binary integrality, and variable
bounds before extracting structured \texttt{BranchRef} outputs.  It does not
claim a full post-reconfiguration AC/DC PF or OPF certificate; callers must read
\texttt{TopoReconfResult::validity.post\_power\_flow\_validated} and
\texttt{full\_hybrid\_opf\_validated}, both currently false for this path.

\subsection{External Solver Integration (HiGHS)}
\label{sec:onr_highs}

The multi-period topology reconfiguration schedule
(\texttt{solve\_topology\_reconfiguration\_schedule()}) supports
automatic dispatch to external MILP solvers.  At runtime the solver selection
proceeds as follows:

\begin{enumerate}
  \item If \texttt{enforce\_ac\_radiality = true} (default) and
        \texttt{tr\_solver = Auto}, the engine checks for
        \texttt{GurobiAdapter}, then \texttt{HighsAdapter} availability.
  \item If an external solver is found, the full multi-period MILP (spanning-tree
        constraints, LinDistFlow, ESS scheduling, switching-cost limits) is
        built and solved via \texttt{solve\_multiperiod\_onr\_highs()}, producing
        $5m + n$ decision variables per timestep.
  \item If no external solver is available, the engine falls back to a
        stepwise per-timestep solve using the native branch-and-cut engine
        (\texttt{solve\_optimal\_reconfiguration()}).
\end{enumerate}

\paragraph{HiGHS Adapter Architecture.}
The \texttt{HighsAdapter} invokes the HiGHS CLI executable via the file-based
protocol:
\begin{enumerate}
  \item \texttt{write\_lp\_as\_mps()} serialises the \texttt{MIPModel} to a
        temporary MPS file with correct \texttt{INTORG}/\texttt{INTEND} markers
        for integer variable blocks.
  \item The adapter constructs a shell command:
\begin{verbatim}
highs --model_file <tmp>.mps --solution_file <tmp>.sol
\end{verbatim}
  \item The solution file is parsed via \texttt{parse\_highs\_model\_status()}
        to extract the model status (Optimal/Feasible/Infeasible) and primal
        variable values.
\end{enumerate}

\paragraph{Executable Resolution.}
The HiGHS binary is resolved through a three-level fallback:
\begin{enumerate}
  \item Explicit path passed to the adapter constructor.
  \item Environment variable \texttt{HACDCPF\_HIGHS\_EXECUTABLE}.
  \item CMake compile-time define \texttt{HACDCPF\_HIGHS\_EXECUTABLE} (set
        by \texttt{find\_program()} in \texttt{Dependencies.cmake}).
\end{enumerate}

\subsection{Current Modeling Boundary}

The ONR module is already useful for radial feeder reconfiguration, but its
current mathematical scope is intentionally narrower than the main PF/OPF stack.
The main structural limitations are:
\begin{itemize}
  \item no hybrid-side injections from \texttt{VPP}, \texttt{Microgrid}, or \texttt{MobileStorage},
  \item no integer co-optimization of DER dispatch and topology in the steady ONR public API,
  \item fault-triggered hybrid results are graph/LinDistFlow topology certificates,
        not post-solve full AC/DC PF or OPF certificates; the result
        \texttt{model\_scope} and validity flags expose this boundary.
\end{itemize}
